\rhead{CY-EL1: Digital Forensics \& Incident Response}
\phantomsection 
\section*{Digital Forensics \& Incident Response (CY-EL1)}
\label{elec:CY_1}

\noindent \small{\hyperref[main:scheme]{$\leftarrow$ Back to Scheme}} \vspace{0.5cm}

\noindent \textbf{Credits:} 4 (3-0-2)

\subsection*{Theory}
\noindent\textbf{Unit 1} \\
\textit{Digital Forensics Fundamentals:} Forensic science principles (Locard's exchange principle, scientific method), Digital evidence types and characteristics (volatile, non-volatile, latent), Chain of custody requirements, Legal considerations (admissibility, Daubert standard), Incident response lifecycle (NIST, SANS models), Forensic readiness planning, Roles (first responders, investigators, analysts).

\vspace{0.3cm}

\noindent\textbf{Unit 2} \\
\textit{Incident Response Processes:} Preparation phase (policies, tools, teams), Detection and analysis (indicators of compromise, triage), Containment strategies (short-term, long-term), Eradication and recovery procedures, Post-incident activities (lessons learned, reporting), Incident classification (levels, severity), Communication protocols and stakeholder management.

\vspace{0.3cm}

\noindent\textbf{Unit 3} \\
\textit{Disk and File System Forensics:} Acquisition methods (dead disk: dd, FTK Imager; live response), Write-blockers and validation (hashing: MD5, SHA-256), File system analysis (FAT, NTFS, ext4), Timeline reconstruction (MFT, \$LogFile, superblocks), File carving (scalpel, foremost), Slack space and unallocated clusters, Metadata extraction (EXIF, MAC times).

\vspace{0.3cm}

\noindent\textbf{Unit 4} \\
\textit{Memory and Application Forensics:} Memory acquisition (AVML, DumpIt, LiME), Volatility framework (imageinfo, pslist, psscan, hashdump), Process injection detection (hollow process, DLL injection), Registry analysis (userassist, amcache, shimcache), Browser forensics (SQLite artifacts: history, downloads, cache), Email forensics (Outlook PST/OST, MBOX), Anti-forensics techniques (timestomping, overwriting).

\vspace{0.3cm}

\noindent\textbf{Unit 5} \\
\textit{Network and Advanced Forensics:} Packet capture analysis (Wireshark, tcpdump), Network forensics (IoCs: C2 beacons, anomalies), Log analysis (Windows Event Logs, Sysmon, firewall logs), Cloud forensics (AWS/Azure artifacts), Mobile forensics (Android logical/physical, iOS backups), Malware forensics (static/dynamic analysis), SIEM integration (Splunk, ELK Stack).

\newpage

\subsection*{Practical}
\noindent\textbf{Lab 1: Forensic Acquisition and Hashing} \\
Focuses on disk imaging with FTK Imager/dd, chain of custody forms, hash verification.

\vspace{0.2cm}

\noindent\textbf{Lab 2: File System Analysis} \\
Focuses on Autopsy timeline reconstruction, file carving, metadata extraction.

\vspace{0.2cm}

\noindent\textbf{Lab 3: Incident Response Simulation} \\
Focuses on NIST IR lifecycle, IOC identification, containment planning.

\vspace{0.2cm}

\noindent\textbf{Lab 4: Memory Forensics with Volatility} \\
Focuses on process analysis, network connections, registry hive extraction.

\vspace{0.2cm}

\noindent\textbf{Lab 5: Registry and Application Artifacts} \\
Focuses on userassist, recent documents, browser history extraction.

\vspace{0.2cm}

\noindent\textbf{Lab 6: Network Forensics} \\
Focuses on Wireshark packet analysis, protocol dissection, anomaly detection.

\vspace{0.2cm}

\noindent\textbf{Lab 7: Log Analysis and Correlation} \\
Focuses on Sysmon/Event Logs parsing, timeline correlation.

\vspace{0.2cm}

\noindent\textbf{Lab 8: Mobile Device Forensics} \\
Focuses on ADB backup, Cellebrite/ Oxygen extraction, app data analysis.

\vspace{0.2cm}

\noindent\textbf{Lab 9: Malware Reverse Engineering Basics} \\
Focuses on strings analysis, YARA rules, sandbox detonation.

\vspace{0.2cm}

\noindent\textbf{Lab 10: Complete Incident Investigation} \\
Focuses on end-to-end forensics report generation, court-ready evidence package.

\vspace{0.2cm}

\newpage
\rhead{Curriculum Schema}
